\documentclass[12pt]{article}
% Préambule commun à tous les documents extraits de « La Revue CAPES »
\usepackage[T1]{fontenc}
\usepackage[utf8]{inputenc}
\usepackage{lmodern}
\usepackage{amsmath,amssymb,amsthm,amsfonts,mathrsfs}
\usepackage{setspace}
\usepackage{listings}
\usepackage{pgfplots}
\pgfplotsset{compat=1.18}
\usepackage{longtable}
\usepackage{adjustbox}
\usepackage{lettrine}
\usepackage{tkz-tab}
\usepackage{changepage}
\usepackage{pdflscape}
\usepackage{graphicx}
\usepackage{tikz}
\usepackage{xcolor}
\usepackage[a4paper,top=2cm,bottom=2.5cm,left=2.5cm,right=2cm]{geometry}
\usepackage{fancyhdr}
\usepackage{draftwatermark}
\SetWatermarkText{La RevueCapes}
\SetWatermarkLightness{0.9}
\SetWatermarkScale{2}
\usepackage{hyperref}
\hypersetup{colorlinks=true,linkcolor=blue!50!black,urlcolor=blue!50!black}

\definecolor{revue}{RGB}{20,60,120}

\pagestyle{fancy}
\fancyhf{}
\renewcommand{\headrulewidth}{0.4pt}
\setlength{\headheight}{15pt}

% \entete{Matière}{Titre}{Type de document}
\newcommand{\entete}[3]{%
  \fancyhead[L]{\small\textsc{La Revue CAPES} -- #1}%
  \fancyhead[R]{\small #2}%
  \fancyfoot[C]{\thepage}%
  \thispagestyle{plain}%
  \begin{center}
    {\color{revue}\rule{\textwidth}{1.2pt}}\\[4pt]
    {\small\textsc{Concours directs CAP-CEG / CAPES -- Burkina Faso}}\\[6pt]
    {\LARGE\bfseries\color{revue} #2}\\[6pt]
    {\large\itshape #1 -- #3}\\[2pt]
    {\color{revue}\rule{\textwidth}{1.2pt}}
  \end{center}
  \vspace{0.5cm}%
}

% Encadré pour les remarques ajoutées dans les corrigés
\newcommand{\remarque}[1]{\par\medskip\noindent\fcolorbox{revue}{revue!6}{\parbox{\dimexpr\linewidth-2\fboxsep-2\fboxrule}{\textbf{\color{revue}Remarque.} #1}}\par\medskip}


\begin{document}
\entete{Culture générale}{Corrigé du sujet type de dissertation n°11}{Correction}
\begin{spacing}{1.5}

\noindent\textbf{Sujet.} L'éducation des jeunes filles au Burkina Faso : enjeux et perspectives.

\rubrique{1. Analyse du sujet}

\begin{itemize}
\item \textbf{Type de sujet :} sujet-thème appelant un \textbf{plan analytique} : I. les enjeux (pourquoi éduquer les filles), II. les obstacles actuels, III. les perspectives (solutions et avenir).
\item \textbf{Mots-clés :}
  \begin{itemize}
  \item « éducation des jeunes filles » : scolarisation, maintien à l'école, réussite scolaire et formation des filles, mais aussi éducation non formelle et alphabétisation ;
  \item « enjeux » : ce qui est en jeu pour les filles elles-mêmes, pour les familles et pour le pays ;
  \item « perspectives » : orientations, solutions et évolutions possibles.
  \end{itemize}
\item \textbf{Problématique :} pourquoi l'éducation des filles est-elle décisive pour le développement du Burkina Faso, et comment lever les obstacles qui l'entravent encore ?
\end{itemize}

\rubrique{2. Plan détaillé}

\begin{enumerate}
\item[I.] \textbf{Les enjeux de l'éducation des filles.} 1. Un enjeu d'équité et de droits. 2. Un enjeu de santé et de bien-être familial. 3. Un enjeu de développement économique et social.
\item[II.] \textbf{Les obstacles persistants.} 1. Les pesanteurs socioculturelles (mariages précoces, travaux domestiques, préjugés). 2. La pauvreté et les insuffisances de l'offre scolaire. 3. L'insécurité et les violences.
\item[III.] \textbf{Les perspectives.} 1. Sensibiliser et faire évoluer les mentalités. 2. Adapter l'école aux besoins des filles. 3. Soutenir les familles et valoriser les modèles de réussite.
\end{enumerate}

\rubrique{3. Proposition de rédaction}

\partie{Introduction}

« Éduquer un garçon, c'est éduquer un individu ; éduquer une fille, c'est éduquer une nation. » Ce proverbe, souvent cité en Afrique, souligne le rôle décisif des femmes dans la famille et dans la société. Au Burkina Faso, la scolarisation des filles a beaucoup progressé au cours des dernières décennies, notamment dans l'enseignement primaire, grâce à l'action de l'État, des partenaires et des associations. Pourtant, les filles restent plus nombreuses que les garçons à abandonner l'école avant la fin du post-primaire et du secondaire, et elles sont moins présentes dans l'enseignement supérieur et les filières scientifiques. Quels sont les enjeux de l'éducation des jeunes filles au Burkina Faso, et quelles perspectives peut-on envisager pour l'améliorer ? Nous analyserons d'abord les enjeux de l'éducation des filles, puis les obstacles qui la freinent, avant de dégager des perspectives.

\partie{I. Les enjeux de l'éducation des filles}

\souspartie{1. Un enjeu d'équité et de droits}
Le droit à l'éducation appartient à tous les enfants, sans distinction de sexe. Priver une fille d'école, c'est la condamner souvent à la dépendance et à la pauvreté. L'éducation donne aux filles les moyens de connaître leurs droits, de prendre la parole, de participer aux décisions et d'accéder à des responsabilités dans la société.

\souspartie{2. Un enjeu de santé et de bien-être familial}
Une femme instruite se marie plus tard, a des enfants plus espacés, connaît mieux les règles d'hygiène, de nutrition et de prévention des maladies, et fait vacciner ses enfants. L'éducation des mères est l'un des facteurs les plus efficaces de réduction de la mortalité maternelle et infantile. Une mère instruite veille aussi davantage à la scolarisation de ses propres enfants.

\souspartie{3. Un enjeu de développement économique et social}
Les femmes jouent un rôle essentiel dans l'agriculture, le commerce et la transformation des produits locaux. Instruites et formées, elles augmentent leurs revenus, créent des entreprises et contribuent davantage à la croissance du pays. Exclure les filles de l'éducation, c'est se priver de la moitié des talents de la nation.

\transition{Malgré ces enjeux évidents, l'éducation des filles se heurte encore à de nombreux obstacles.}

\partie{II. Les obstacles persistants}

\souspartie{1. Les pesanteurs socioculturelles}
Dans certaines familles, la scolarisation des filles est jugée moins utile que celle des garçons, car la fille est destinée à se marier et à quitter la famille. Les mariages précoces ou forcés interrompent la scolarité de nombreuses adolescentes. Les filles sont aussi davantage sollicitées pour les travaux domestiques (corvée d'eau, garde des petits frères, cuisine), ce qui réduit leur temps d'étude.

\souspartie{2. La pauvreté et l'insuffisance de l'offre scolaire}
Lorsque les ressources sont limitées, les familles privilégient souvent la scolarisation des garçons. L'éloignement des établissements, notamment au post-primaire, expose les filles à des risques sur le trajet ou les oblige à vivre loin de leur famille. L'absence de latrines séparées et de points d'eau dans certaines écoles pénalise particulièrement les adolescentes.

\souspartie{3. L'insécurité et les violences}
Les grossesses précoces, parfois liées à des violences sexuelles, y compris en milieu scolaire, entraînent de nombreux abandons. Par ailleurs, la crise sécuritaire a fermé de nombreuses écoles et rendu les filles déplacées plus vulnérables aux mariages précoces et à l'exploitation.

\transition{Ces obstacles ne sont pas insurmontables : plusieurs perspectives permettent d'améliorer l'éducation des filles.}

\partie{III. Les perspectives}

\souspartie{1. Sensibiliser et faire évoluer les mentalités}
Les campagnes de sensibilisation menées auprès des parents, des leaders coutumiers et religieux, ainsi que l'action des associations de mères éducatrices, permettent de convaincre les communautés de l'importance de l'éducation des filles. L'application effective de la loi contre les mariages d'enfants est indispensable.

\souspartie{2. Adapter l'école aux besoins des filles}
Il faut rapprocher les établissements des lieux d'habitation, construire des latrines séparées, développer les cantines et les internats sécurisés, lutter fermement contre le harcèlement et les violences en milieu scolaire et proposer une éducation à la santé sexuelle adaptée. Les filières techniques et scientifiques doivent être ouvertes et rendues attractives pour les filles.

\souspartie{3. Soutenir les familles et valoriser les modèles de réussite}
Les bourses, les kits scolaires et les aides aux familles vulnérables réduisent le coût de la scolarisation. La mise en avant de femmes ayant réussi — enseignantes, médecins, ingénieures, entrepreneures, responsables politiques — encourage les filles à poursuivre leurs études. L'alphabétisation des femmes adultes et la formation professionnelle offrent une seconde chance à celles qui ont quitté l'école.

\partie{Conclusion}

L'éducation des jeunes filles est un enjeu majeur pour le Burkina Faso : elle est une question de justice, une condition de la santé des familles et un moteur du développement économique et social. Elle se heurte cependant à des obstacles socioculturels, économiques et sécuritaires qui provoquent encore de nombreux abandons. En sensibilisant les communautés, en adaptant l'école aux besoins des filles et en soutenant les familles, il est possible de relever ce défi. Investir dans l'éducation des filles, c'est investir dans l'avenir de toute la nation.

\end{spacing}
\end{document}
